% Source note: notes/2026-11-02-mt4.md
\begin{entry}{research}{MT4}{Research, product rethink, simulator plan and setup}
\entrytags{users, aad, synthetic-eeg}

\subsection*{Who needs this, and can it work?}
\entryfield{What I did}
Read papers on (1) how people with hearing loss experience music, and (2) whether you can tell from EEG which instrument someone is listening to.
\entryfield{What I found}
\begin{itemize}
  \item Hearing aids are bad for music. In a survey of 176 hearing aid users, 67\% had some trouble with music \cite{Greasley2020}.
  \item Cochlear implant users liked music better when the vocals were turned up by about 8\,dB. People with normal hearing preferred the original \cite{Tahmasebi2020}.
  \item Telling which instrument someone is listening to from EEG works, but not very well: about 71\% with a full lab cap over 60\,s \cite{Simon2023}, and about 60\% for implant users \cite{Althoff2026}.
  \item A Sony study got about 86\% using a cheap 4-electrode headband (Muse 2), but only tested 8 people \cite{Akama2026}.
  \item Brain-reading models trained on speech don't work on music \cite{Simon2023}.
  \item Headsets with electrodes only around the ear do much worse than a full cap \cite{Zhu2024}. Real-time systems take 5--12\,s to notice when you switch attention \cite{COCOHA2024}.
\end{itemize}
\entryfield{What I learned}
\begin{itemize}
  \item There is an identifiable target: people with hearing loss who struggle with music.
  \item A simple fixed setting (``vocals always louder'') already helps them, so our product has to beat that.
  \item There is very little EEG data recorded during music, which is a reason to generate fake data (my area).
\end{itemize}

\subsection*{Rethinking what the product is for}
\entryfield{What I did}
Questioned whether the MT3 idea made sense, and who it is actually for.
\entryfield{The problem with the MT3 pitch}
\begin{itemize}
  \item For someone with normal hearing, a volume slider does the job better than EEG. It's faster, cheaper and more reliable.
  \item We listed too many possible users.
\end{itemize}
\entryfield{New version}
\begin{itemize}
  \item Users: people with hearing aids or cochlear implants, who find music muddy.
  \item Why turning it all up doesn't help: hearing loss blurs sounds together, and louder doesn't separate them.
  \item Product: a music player that slowly turns up the part you're following and turns the rest down a little. It only changes when it's confident, and otherwise stays on your normal setting.
  \item Why EEG and not a button: it keeps up as your attention moves during a song, and works when your hands or eyes are busy.
\end{itemize}
\entryfield{What changes for testing}
Compare against a fixed setting like ``vocals louder'', not just against the original song.

\subsection*{Planning the EEG simulator}
\entryfield{What I did}
Worked out what the simulator should produce, checked whether something like it already exists, and found a dataset to build it from.
\entryfield{What it makes}
Fake EEG of someone listening to a song while focusing on one instrument, as it would look on a given headset. Built in steps:
\begin{enumerate}
  \item How loud each instrument is over time.
  \item How the brain responds to that loudness, learned from real recordings.
  \item A boost for the instrument being listened to.
  \item Background brain activity, taken from real recordings.
  \item Noise like blinks, and the electrode positions of a chosen headset.
\end{enumerate}
\entryfield{What it's for}
Training data (real music EEG is extremely limited), and predicting how well a cheap headset would work before buying one.
\entryfield{Does it already exist?}
Tools exist for parts of it: simulating EEG in general \cite{Gramfort2013}, background brain activity \cite{Cole2019}, fitting brain responses \cite{Brodbeck2023}, and simulating EEG for short events in MATLAB \cite{Krol2018}. One group generated fake EEG with AI for speech, and it barely helped \cite{Rannaleet2026}. I couldn't find anything that turns music into labelled fake EEG for a given headset. So what's new is the combination and the experiments, not the individual parts.
\entryfield{Dataset}
MAD-EEG \cite{Cantisani2019}. Downloaded it. It has 6 usable people and no electrodes near the ears, so it can't directly tell us about ear-based headsets.
\entryfield{Plan for MT5}
Check the dataset shows a clear brain response to the music; get a baseline accuracy on real data with Cody; build a first simple version of the simulator; get the model of the headset we will decide on from the team / Maxime's research.
\subsection*{Setting up the repo and data}
\entryfield{What I did}
\begin{itemize}
  \item Created the team GitHub repo: \url{https://github.com/b-adamson/attention-earbuds}. Kept it minimal (README, data notes) so the team can agree the structure together.
  \item Downloaded MAD-EEG \cite{Cantisani2019} and added a small loader (\texttt{simulation/load\_madeeg.py}) that reads one trial into a Python dictionary.
\end{itemize}
\entryfield{What the data looks like}
\begin{itemize}
  \item Stored as HDF5: one file organised like folders. Subject $\rightarrow$ trial $\rightarrow$ arrays. A YAML file describes each trial (instruments, which one was attended, sample rates).
  \item Each trial has the EEG (20 electrodes, 256\,Hz), the stereo mix heard, and each instrument's track on its own (44.1\,kHz). Duos have an empty third instrument slot.
  \item 6 subjects and 246 trials in the preprocessed file. No electrodes near the ears.
\end{itemize}
\entryfield{Next}
Write my own loudness extraction and brain-response (TRF) fit, and check for a clear response in the real data.

\subsection*{Team meeting (08/11/2026)}
\entryfield{Discussion doc}
\url{https://docs.google.com/document/d/13DpkpDFB5p7-QVpHNyvnnZCEQ5L7WrmuGmXRaS7t8ks/edit}
\entryfield{My contributions}
\begin{itemize}
  \item Proposed focusing on hearing aid / cochlear implant users, and testing against a fixed ``vocals louder'' setting.
  \item Shared my progress on the simulator and the MAD-EEG dataset.
  \item Suggested a folder per person in the team repo, and a shared data format between the simulator and Cody's decoder.
\end{itemize}
\entryfield{Actions}
\begin{itemize}
  \item Billy: check the music response in MAD-EEG is strong enough to build the simulator on.
  \item Billy \& Cody: agree data format between simulator and decoder.
  \item Maxime: send electrode positions for shortlisted headsets.
\end{itemize}
% TODO(Billy): add team decisions once everyone has filled in the doc.
\end{entry}
